\section{控制分配：通道解耦的几何机理}
\label{sec:allocation}

控制分配（control allocation）是连接控制律输出（期望力矩）与执行器输入（摆角/差速指令）
的数学环节。本构型的正交摆轴使推力力臂主项形成对角结构，因而可用三通道直接映射作为
低复杂度近似；反扭矩投影形成的交叉项仍需评估或补偿。本章全部矩阵先在推进模型系
$\{P\}$中推导，令$\bm M^P=[M_{x'},M_{y'},M_{z'}]^T$；固件任务通道经
式(\ref{eq:frame_permutation})置换后解释。

\subsection{小角度近似与线性化效能矩阵}

在小摆角近似下（$\sin\delta \approx \delta$，$\cos\delta \approx 1$，$\delta \ll 1$），
并假设前后推力近似相等（$T_f \approx T_t \approx T_0$），式(\ref{eq:six_component})
中的力矩方程线性化为：
\begin{equation}
    \begin{bmatrix} M_{x'} \\ M_{y'} \\ M_{z'} \end{bmatrix}
    \approx
    \underbrace{
    \begin{bmatrix}
        -k_Q \omega_f^2 + k_Q \omega_t^2 \\
        -b T_t \delta_t \\
        a T_f \delta_f
    \end{bmatrix}
    }_{\text{主控通道}}
    +
    \underbrace{
    \begin{bmatrix}
        0 \\
        -\tau_f \delta_f \\
        -\tau_t \delta_t
    \end{bmatrix}
    }_{\text{反扭矩投影交叉项}}
    \label{eq:moment_linearized}
\end{equation}

\textbf{差速变量的参数化}：定义差速变量$\Delta\omega$，使两电机目标转速为：
\begin{equation}
    \omega_f^{\text{target}} = \omega_0 \sqrt{1 + \Delta\omega}, \quad
    \omega_t^{\text{target}} = \omega_0 \sqrt{1 - \Delta\omega}
    \label{eq:diff_def}
\end{equation}
其中$\omega_0$为基准转速（由油门决定）。平方根的数学定义域为$[-1,1]$，
当前固件配置限幅为$|\Delta\omega|\leq0.7$。在单侧目标转速尚未触发
$\omega_{\max}$钳位时，$\omega_f^2+\omega_t^2=2\omega_0^2$严格成立；
一旦钳位，该恒等式和总推力保持性质均被破坏。由$T\propto\omega^2$，
一阶近似下总推力保持不变，差速仅影响扭矩差而不影响总推力。
这一性质在未饱和的局部模型中实现油门--差速解耦；经当前固件轴置换后，
差速对应偏航任务通道，而不是滚转通道。

令基准反扭矩$\tau_0=k_Q\omega_0^2$，模型$x'$轴力矩的差速分量为：
\begin{equation}
    \Delta M_{x'}^{\text{差速}} = k_Q (\omega_t^2 - \omega_f^2)
    = k_Q (\omega_0^2(1-\Delta\omega) - \omega_0^2(1+\Delta\omega))
    = -2 k_Q \omega_0^2 \Delta\omega
    = -2 \tau_0 \Delta\omega
    \label{eq:differential_roll}
\end{equation}

在零摆角、等基准转速处，对输入$[\Delta\omega,\delta_t,\delta_f]^T$求完整一阶Jacobian，得到：
\begin{equation}
    \boxed{
    \begin{bmatrix} M_{x'} \\ M_{y'} \\ M_{z'} \end{bmatrix}
    =
    \underbrace{
    \begin{bmatrix}
        -2\tau_0 & 0 & 0 \\
        0 & -b T_0 & -\tau_0 \\
        0 & -\tau_0 & a T_0
    \end{bmatrix}
    }_{\bm{B}_{\mathrm{full}}(\omega_0)}
    \begin{bmatrix} \Delta\omega \\ \delta_t \\ \delta_f \end{bmatrix}
    + \text{气动项}
    }
    \label{eq:effectiveness_matrix}
\end{equation}

若忽略相对较小的$\tau_0$交叉项，则得到正文后续直接映射所用的主通道近似矩阵
$\bm{B}_{\mathrm{main}}=\operatorname{diag}(-2\tau_0,-bT_0,aT_0)$。

\textbf{适用范围说明}：$\mathbf{B}_{\mathrm{full}}$是在$\delta_f=\delta_t=0$处对非线性力矩映射
（式\ref{eq:six_component}）做一阶Taylor展开得到的局部Jacobian，适用于解释主通道结构。
有限摆角时应使用工作点相关矩阵
\begin{equation}
\boxed{
\bm B_{\mathrm{true}}=
\begin{bmatrix}
-k_Q\omega_0^2(c_f+c_t) & -Q_t s_t & Q_f s_f\\
b k_T\omega_0^2s_t-k_Q\omega_0^2s_f & -bT_t c_t & -Q_f c_f\\
a k_T\omega_0^2s_f+k_Q\omega_0^2s_t & -Q_t c_t & aT_f c_f
\end{bmatrix}}
\label{eq:Btrue_allocation}
\end{equation}
其中$s_i=\sin\delta_i$、$c_i=\cos\delta_i$。式(\ref{eq:Btrue_allocation})表明非零摆角时
差速列也会投影到$M_{y'}$和$M_{z'}$。当前固件在线构造该$3\times3$矩阵并计算绝对指令；
Web默认仍采用直接SAS映射，但可通过\texttt{useBtrue}分支启用同类在线矩阵。

\subsection{主通道近似对角化的几何根源}

推力力臂贡献在$\bm{B}_{\mathrm{main}}$中为对角结构，这是正交摆轴与纵向力臂共同作用的结果；
完整矩阵$\bm{B}_{\mathrm{full}}$并非严格对角。主通道结构可由以下几何关系解释：

\begin{itemize}
    \item \textbf{前电机绕$z'$轴摆动}：推力方向的变化限制在$x'y'$平面内。
          推力$y'$分量经力臂$a$产生模型$z'$主力矩$M_{z'}=aT_f\sin\delta_f$。
          前电机推力不产生$z'$向分量，因此其力臂项不产生$M_{y'}$
          ——但反扭矩投影仍产生一阶小角度耦合$-\tau_f\delta_f$；
    \item \textbf{尾电机绕$y'$轴摆动}：推力方向的变化限制在$x'z'$平面内。
          推力$z'$分量经力臂$b$产生模型$y'$主力矩$M_{y'}=-bT_t\sin\delta_t$。
          尾电机推力不产生$y'$向分量，因此其力臂项不产生$M_{z'}$
          ——但反扭矩投影仍产生一阶小角度耦合$-\tau_t\delta_t$；
    \item \textbf{差速通道}：在零摆角名义点只影响$M_{x'}$；非零摆角时，
          转速变化会同时改变反扭矩投影和摆动力矩幅值，故式(\ref{eq:Btrue_allocation})其余两行一般不为零。
\end{itemize}

这一设计使控制分配问题从通用的$\bm{u}=\bm{B}^+\bm{M}_{\text{cmd}}$（若
$\bm B\in\mathbb R^{3\times3n}$，则$\bm B^+\in\mathbb R^{3n\times3}$，$3n>3$）
可近似为简单的\textbf{三通道主通道映射}：
\begin{equation}
    \Delta\omega = -\frac{M_{x'}^{\text{cmd}}}{2\tau_0}, \quad
    \delta_t = -\frac{M_{y'}^{\text{cmd}}}{b T_0}, \quad
    \delta_f = \frac{M_{z'}^{\text{cmd}}}{a T_0}
    \label{eq:direct_mapping}
\end{equation}

式(\ref{eq:direct_mapping})是理论上的低复杂度分配，也是Web默认SAS的主通道近似；
Web可选分支与当前固件则使用在线矩阵。三者应分开陈述。
其适用范围受低油门、执行器饱和和交叉项量级限制\cite{johansen2013control,oppenheimer2010control}。

\subsection{耦合项量级评估与补偿策略}

线性化后的交叉耦合项来源于反扭矩在非主控轴方向的投影。在小摆角下它们对摆角是一阶项：
\begin{align}
    \Delta M_{y'}^{\text{耦合}} &= -\tau_f \delta_f
    \approx -k_Q \omega_f^2 \cdot \delta_f
    \label{eq:coupling_pitch} \\
    \Delta M_{z'}^{\text{耦合}} &= -\tau_t \delta_t
    \approx -k_Q \omega_t^2 \cdot \delta_t
    \label{eq:coupling_yaw}
\end{align}

\textbf{耦合-主控比}的量级评估。以俯仰通道为例，
主控力矩量级为$b T_t \delta_t \approx b \cdot k_T \omega_0^2 \cdot \delta_t$，
耦合力矩量级为$\tau_f \delta_f \approx k_Q \omega_0^2 \cdot \delta_f$。在俯仰和偏航
同时大幅偏转的极限工况，耦合比约为：
\begin{equation}
    \frac{\text{耦合}}{\text{主控}}
    = \frac{k_Q \omega_0^2 \delta_f}{b \cdot k_T \omega_0^2 \delta_t}
    = \frac{k_Q}{k_T b} \cdot \frac{\delta_f}{\delta_t}
    = \frac{C_Q}{C_T} \cdot \frac{D}{b} \cdot \frac{\delta_f}{\delta_t}
    \label{eq:coupling_ratio}
\end{equation}

由量纲分析$k_Q / (k_T b) = (C_Q / C_T) \cdot (D / b)$。
当前MODEL-DEFAULT参数给出$k_Q/(k_Tb)\approx0.08547$（$b=0.315$\,m）。
实机三轮调参只提供控制效果的量级锚点，不等同于$k_T$、$k_Q$独立扭矩台架标定；该比值
只说明交叉项不能被当作严格为零。实际耦合还随两通道摆角比、转速差与饱和状态变化。

\textbf{工程对策}：
\begin{enumerate}[label=(\arabic*)]
    \item \textbf{显式前馈补偿}：在控制律中直接计算耦合项并前馈抵消，例如
          \begin{equation*}
              \Delta \delta_t^{\text{comp}}
              = \frac{k_Q \omega_f^2}{b T_0}\,\delta_f .
          \end{equation*}
          该补偿仍需结合符号约定和执行器饱和进行闭环验证；
    \item \textbf{SAS鲁棒性}：设计足够高的SAS反馈增益，使闭环系统对耦合扰动具有足够衰减。
          耦合扰动在闭环下被抑制了$1/(1+GH)$倍，其中$GH$为回路增益；
    \item \textbf{工作点Jacobian或非线性求解}：在线使用式(\ref{eq:Btrue_allocation})可保留局部交叉项；
          更大有效域可用非线性映射（式\ref{eq:six_component}）配合约束求解。
          $15^\circ$时$\sin\delta/\delta\approx0.989$，单独由正弦近似产生的误差仅约$1.1\%$，
          因而不能据此宣称Newton迭代“必需”；固件当前限制为$\pm15^\circ$，Web为$\pm25^\circ$。
\end{enumerate}

\subsection{低油门差速控制退化与增益调度}

差速效能$\partial M_{x'}/\partial\Delta\omega=-2k_Q\omega_0^2$与基准推力同为
$\omega_0^2$量级；两摆座的力臂效能也随$T_i\propto\omega_i^2$下降。因此低油门并非
只削弱某一个轴，而是同时压缩全部推进姿态控制权限，其中差速模型轴在当前固件任务系中
对应偏航。提高反馈增益不能创造已经消失的物理力矩，且会把未建模滞后和噪声放大。

在当前$\omega_0=\mathrm{thr}\,\omega_{\max}$指令映射下，油门由100\%降到50\%时，
名义效能降到25\%；油门趋近零时，推进系统的三轴姿态效能均趋近零。该平方关系是
当前模型的指令假设，不是跨电调、电池和螺旋桨组合均成立的普适定律。飞行包线设计应据此
设置最小可用推力和控制裕度，而不能直接把固定翼低油门进近场景等同于悬停差速控制问题。

\textbf{增益调度与瞬态失配（2026-08-09项目记录）}可分为三层：
\begin{enumerate}[label=(\arabic*)]
    \item 低油门按$1/w_0^2$归一化差速指令会产生数值放大，调度改用$(w_0/w_h)^2$；
    \item 物理力矩仍与$w_0^2$同阶，稳态调度需封顶以避免高油门增益继续放大；
    \item 指令转速与实际响应受$\tau_m$滞后分离，油门瞬态会造成分配工作点失配。
\end{enumerate}
固件采用开环模型观测器
$w_{\mathrm{est}}\mathrel{+}=(w_{\mathrm{cmd}}-w_{\mathrm{est}})\Delta t/\tau_m$缓解该失配；
它并非测速传感器。当前实现还具有不对称工作点：$\bm B_{\mathrm{true}}$的差速列使用输入结构中的
指令基准转速$w_0$，摆角列中的$T_f,T_t,Q_f,Q_t$使用经下限处理的观测状态转速。
因此不能概括为“B矩阵全部使用同一个观测转速”。固件最终还实施
$|\delta_f|,|\delta_t|\le15^\circ$、$|\Delta\omega|\le0.7$以及5\%零油门归中门控。

\subsection{与通用过驱动推力矢量平台的对比}

通用多推进器矢量推力平台的设计空间可以是过驱动（over-actuated）或全驱动（fully-actuated）。
对于$n$个推进器、每个推进器具有2自由度万向节的情况，执行空间维数为$3n$
（$n$个推进器$\times$ 3维力矢量），映射到6维力/力矩空间。
当$3n > 6$（即$n \ge 3$）时形成过驱动系统\cite{ding2018modeling}。

过驱动系统的控制分配是一个经典的欠定问题：
给定期望力矩$\bm{M}_{\text{cmd}} \in \mathbb{R}^3$（或期望六维力/力矩$\in \mathbb{R}^6$），
求执行器输入$\bm{u} \in \mathbb{R}^{3n}$，使$\bm{B}\bm{u} = \bm{M}_{\text{cmd}}$。
由于$3n > 3$（或$>6$），方程有无穷多解。需要在解空间中选取“最优”解，
通常通过求解以下优化问题之一\cite{johansen2013control}：
\begin{itemize}
    \item \textbf{最小二范数（伪逆）}：$\min_{\bm{u}} \|\bm{u}\|^2$ s.t. $\bm{B}\bm{u} = \bm{M}_{\text{cmd}}$。
          解为$\bm{u} = \bm{B}^T(\bm{B}\bm{B}^T)^{-1}\bm{M}_{\text{cmd}}$。简单但未考虑执行器限幅；
    \item \textbf{加权伪逆}：$\min_{\bm{u}} \bm{u}^T\bm{W}\bm{u}$ s.t. $\bm{B}\bm{u} = \bm{M}_{\text{cmd}}$。
          通过权重矩阵$\bm{W}$引入不同执行器优先级（如优先使用大推力电机、保护接近饱和的执行器）；
    \item \textbf{二次规划}：$\min_{\bm{u}} \|\bm{W}_u(\bm{u} - \bm{u}_p)\|^2$ s.t.
          $\bm{B}\bm{u} = \bm{M}_{\text{cmd}}$且$\bm{u}_{\min} \le \bm{u} \le \bm{u}_{\max}$。
          在满足等式和不等式约束的前提下最小化与标称输入的偏差。
\end{itemize}

\sloppy
在线求解带约束优化问题通常需要active-set或内点等迭代算法。
本构型把姿态分配维数降到三个输入对应三个力矩，使$\bm B$成为固定规模方阵；
矩阵并非严格对角，当前固件仍在线计算稠密$3\times3$矩阵并用Cramer形式求逆。
其工程价值是固定$O(1)$计算量和清晰的主通道结构，而不是“零计算开销”或全局完全解耦。
\fussy
